\subsection{An Existence Proof for the Dirichlet Problem}\label{sec:dirichlet-existence-proof}

In this section, we shall establish existence of solutions for the Dirichlet problem. Specifically, we shall prove the following theorem:

\begin{Theorem}
\label{thm:dirichlet-continuous-data-existence}
Let $\Omega$ be a bounded domain in $\mathbb{R}^n$ with a $C^2$-boundary. Then, for any function $g\in C(\partial\Omega)$, there is a unique $u\in C^2(\Omega)\cap C(\overline\Omega)$ satisfying $\Delta u=0$ in $\Omega$ and $u=g$ on $\partial\Omega$.
\end{Theorem}

It will be evident from the proof that the assumption that $\partial\Omega$ is of class $C^2$ can be relaxed; for example, all convex domains are permissible.

The proof will be based on the ideas of Perron, which make use of the following notions. We call $v$ a subsolution (supersolution) if $\Delta v\geq0$ ($\Delta v\leq0$) in $\Omega$ and $v\leq g$ ($v\geq g$) on $\partial\Omega$. Obviously subsolutions exist, e.g., every sufficiently small constant is a subsolution. The maximum principle (weak form) shows that if $u$ is a solution, then $v\leq u$ for every subsolution (this is why we call them subsolutions). Thus, the pointwise supremum of all subsolutions (which is a well defined function) gives us an obvious candidate for a solution. If we can show that this function actually solves the problem, our existence proof will be complete. Before we prove this, we first need to develop a number of prerequisites.

\subsubsection{The Dirichlet Problem on a Ball}\label{subsec:dirichlet-problem-on-ball}

The most old-fashioned way to prove existence of solutions is to give a formula for them. (We discussed some elementary methods for doing this in Section~\ref{subsec:laplace-equation}.) We now obtain an explicit solution for the Dirichlet problem on a ball.

\begin{Theorem}
\label{thm:poisson-ball-formula}
Let $B$ be the ball of radius $R$ centered at the origin, and let $g$ be continuous on $\partial B$. Then the function
\begin{equation}\label{eq:poisson-ball-solution-formula}
u(\bm{x})=\dfrac{R^2-|\bm{x}|^2}{n\omega_nR}\int_{\partial B}\dfrac{g(\bm{y})}{|\bm{x}-\bm{y}|^n}\,\mathrm{d}S_{\bm{y}}
\end{equation}
is of class $C^2(B)$ and satisfies $\Delta u=0$. Moreover, for every $\bm{y}\in\partial B$, we have
\begin{equation}\label{eq:poisson-ball-boundary-limit}
\lim_{\bm{x}\to\bm{y}}u(\bm{x})=g(\bm{y}).
\end{equation}
Here $\omega_n$ denotes the volume of the unit ball in $\mathbb{R}^n$, and the notation $\mathrm{d}S_{\bm{y}}$ in the surface integral indicates the variable of integration.
\end{Theorem}

Equation \eqref{eq:poisson-ball-solution-formula} is known as Poisson\textquotesingle s formula. We note that the special case $\bm{x}=0$ leads to the well known mean-value property: The value of a harmonic function at a point is equal to its average on any ball centered at that point.

\begin{Proof}
Since we can differentiate under the integral, $u$ is in fact of class $C^\infty$ in $B$ and it is a simple calculation to show that it is harmonic. It remains to establish \eqref{eq:poisson-ball-boundary-limit}. Let
\begin{equation}\label{eq:poisson-ball-kernel}
K(\bm{x},\bm{y})=\dfrac{R^2-|\bm{x}|^2}{n\omega_nR|\bm{x}-\bm{y}|^n},\quad\bm{x}\in B,\ \bm{y}\in\partial B,
\end{equation}
and
\begin{equation}\label{eq:poisson-ball-kernel-normalization-function}
\psi(\bm{x})=\int_{\partial B}K(\bm{x},\bm{y})\,\mathrm{d}S_{\bm{y}}.
\end{equation}
Obviously, \eqref{eq:poisson-ball-kernel-normalization-function} is just the special case $g=1$ in \eqref{eq:poisson-ball-solution-formula}; hence $\psi$ is harmonic. It is also obvious that $\psi$ is a radially symmetric function. For radially symmetric functions, Laplace\textquotesingle s equation reads
\begin{equation}\label{eq:radial-laplace-equation}
u_{rr}+\dfrac{n-1}{r}u_r=0,
\end{equation}
and the only solutions of this equation which are regular at the origin are constants. Hence $\psi(\bm{x})=\psi(0)=1$.

Now let $\bm{x}_0\in\partial B$ and let $\epsilon>0$ be given. Choose $\delta>0$ so that $|g(\bm{y})-g(\bm{x}_0)|<\epsilon$ for $|\bm{y}-\bm{x}_0|<\delta$ and let $M$ be an upper bound for $g$ on $\partial B$. For $|\bm{x}-\bm{x}_0|<\delta/2$, we have
\begin{equation}\label{eq:poisson-ball-boundary-estimate}
\begin{aligned}
|u(\bm{x})-g(\bm{x}_0)|
&=\left|\int_{\partial B}K(\bm{x},\bm{y})(g(\bm{y})-g(\bm{x}_0))\,\mathrm{d}S_{\bm{y}}\right|\\
&\leq\int_{|\bm{y}-\bm{x}_0|\leq\delta}K(\bm{x},\bm{y})|g(\bm{y})-g(\bm{x}_0)|\,\mathrm{d}S_{\bm{y}}\\
&\quad+\int_{|\bm{y}-\bm{x}_0|\geq\delta}K(\bm{x},\bm{y})|g(\bm{y})-g(\bm{x}_0)|\,\mathrm{d}S_{\bm{y}}\\
&\leq\epsilon+\dfrac{2M(R^2-|\bm{x}|^2)R^{n-2}}{(\delta/2)^n}.
\end{aligned}
\end{equation}
As $\bm{x}\to\bm{x}_0$, the last term on the right-hand side tends to zero and the theorem follows.
\end{Proof}
